\subsection{Growth ranges and compact examples}
\label{sec:optimal-growing-jet-ranges}

The quadrature boundary differs from the model circle by $O(R^2)$ in physical coordinates. Comparing this displacement with the transition width $R/\sqrt q$ gives the universal range. A holomorphic change of coordinate produces a compact example attaining this displacement.

\Needspace{24\baselineskip}
\begin{theorem}\label{thm:optimal-growing-jet-range}\index{Bergman kernel!optimal growing-jet range}
Let $z_m$ denote the coordinate values of the observation points. For each fixed integer $r\geq1$ and real $\tau$, the law
\begin{equation}\label{eq:optimal-growing-jet-law}
 \frac{\pi B_{m,q}(z_m)}{mr^2|z_m|^{2r-2}}
 \longrightarrow\Phi(\tau),\qquad
 m|z_m|^{2r}=q/r+\tau\sqrt{q/r},
\end{equation}
holds for every choice of the preceding geometric data and every sequence of observation
directions when
$q=o(m^{1/(r+1)})$. This is the optimal universal range: for every
excluded sequence $q\to\infty$, $q=o(m)$, a compact example satisfying
the same hypotheses violates the law along a subsequence.


Write $\Gamma_t$ as $|z|=R c_R(\theta)$. For a fixed metric, the exact
criterion for \eqref{eq:optimal-growing-jet-law} for every sequence of directions and every
fixed $\tau$ is
\begin{equation}\label{eq:fixed-weight-jet-criterion}
 \sqrt q\,\|c_R-1\|_\infty\longrightarrow0.
\end{equation}
\end{theorem}
\begin{proof}
By \cref{lem:curvature-quadrature-domain}, $c_R-1=O(R)$ in every fixed norm required
here. The observation in \eqref{eq:optimal-growing-jet-law} satisfies
\[
 |z_m|/R=1+\tau/(2\sqrt{rq})+O(q^{-1}).
\]
Since the boundary slope tends to zero, its corrected normal argument
is $\tau+o(1)$ whenever $\sqrt q R\to0$. This condition is precisely
$q=o(m^{1/(r+1)})$. Moreover $a(z)=1+o(1)$ and
$\rho(z)=r^2|z|^{2r-2}(1+o(1))$ there, so the two normalizations agree.
The edge theorem proves universal sufficiency.

The same reasoning proves the sufficiency of \eqref{eq:fixed-weight-jet-criterion}.
For necessity choose, on a subsequence, an angle realizing the maximum
boundary displacement and observe at $|z|=R$ with $\tau=0$. A nonzero
finite subsequential scaled displacement gives a nonzero Gaussian
argument; an infinite one gives saturation to zero or one. Radial and
normal distances have relative error tending to zero, since the graph
has vanishing slope. In either case the limit differs from $1/2$.

We give a compact model for sharpness without assumptions at infinity.
On $\mathbb P^1$ choose a radial metric on $\mathcal O(d)$ with local potential
$f(u)$, $u=|w|^2$, equal to $u^r$ near zero, and with
$g(u)=uf'(u)$ strictly increasing on $(0,\infty)$ and tending to $d$.
For large $u$ require $g(u)=d-\epsilon/u$, with $\epsilon>0$.
Such a smooth interpolation exists for a sufficiently large integer
$d$: prescribe a positive $g'$ equal to $r^2u^{r-1}$ near zero and
$\epsilon/u^2$ near infinity, with total integral $d$.
Then $f=d\log u+\mathrm{constant}+\epsilon/u$ near infinity; the
infinity-frame potential is smooth with strictly positive curvature.
Choose a smooth radial volume equal to $\mathrm{d}A_w$ near zero.

For sufficiently small $t$, the radial hole lies in the exact-model
region and is $|w|<R$: rotational invariance and uniqueness make it a
disk, and its curvature mass is $rR^{2r}=t$. The corrected transition
and saturation theorem therefore apply to this compact example.

Change the allowed observation coordinate to $w=z+c z^2$, $c>0$.
The same compact metric now has
$\phi(z)=|z|^{2r}+O(|z|^{2r+1})$, and the volume still has $a(0)=1$.
Vanishing order is invariant. At the raw observation $z=R$,
\[
 2\sqrt{mr^2R^{2r-2}}\,(|w|-R)
 =2c\sqrt r\,\sqrt q R.
\]
The two curvature-model denominators have ratio tending to one.
Every excluded sequence has a subsequence on which $\sqrt q R$ tends
to a positive finite number or to infinity. \Cref{thm:curvature-boundary-jets}, using its saturation conclusion when the
argument diverges, then gives a value strictly larger than $1/2$. This proves sharpness.
\end{proof}

\subsection{Corrections from the smooth weight}
\label{sec:growing-jet-corrections}
\index{Bergman kernel!effective Taylor correction}

Write
\[
 \phi(z)=|z|^{2r}+P_{2r+1}(z,\bar z)+O(|z|^{2r+2}),\qquad
 P(\theta)=P_{2r+1}(\mathrm{e}^{\mathrm{i}\theta},\mathrm{e}^{-\mathrm{i}\theta}),
\]
where $P_{2r+1}$ is a real homogeneous polynomial. Let $\Lambda$ be the
circle Dirichlet-to-Neumann operator,
$\Lambda \mathrm{e}^{\mathrm{i} n\theta}=|n|\mathrm{e}^{\mathrm{i} n\theta}$. The first boundary coefficient is
\begin{equation}\label{eq:first-boundary-coefficient}
 c_R=1+Rb_1+O(R^2),\qquad
 b_1=-\frac{((2r+1)-\Lambda)P}{4r^2}.
\end{equation}
To derive it, write the interior obstacle to first order as
$1+2r\log s+R H$, with harmonic $H$. Value matching on
$s=1+Rb_1$ gives $H|_{s=1}=P$; the terms involving $b_1$ cancel since
the radial reference already matches first derivatives. Derivative
matching gives $4r^2b_1+(2r+1)P-\Lambda P=0$.
Pure harmonic terms disappear, as required by invariance under frame
changes.

For $q=O(m^{1/(r+1)})$, at the raw observation the corrected argument is
\begin{equation}\label{eq:first-growing-jet-shift}
 \tau+\frac{\sqrt q R}{2r^{3/2}}((2r+1)-\Lambda)P(\theta)+o(1),
\end{equation}
uniformly in $\theta$. The kernel differs from $\Phi$ of the displayed
argument by $o(1)$. The error is small because $\sqrt q R^2\to0$ in
this range. For $w=z+cz^2$, $P=2r\operatorname{Re}(c\mathrm{e}^{\mathrm{i}\theta})$ and the shift is
$2\sqrt{rq}R\operatorname{Re}(c\mathrm{e}^{\mathrm{i}\theta})$, agreeing with the compact example.

For every fixed integer $J\geq1$ the parameter construction gives
\[
 c_R=1+\sum_{j=1}^J R^j b_j+O(R^{J+1})
\]
in every fixed norm needed. There is an explicit recursion: expand the
interior harmonic obstacle in powers of $R$. At order $j$, value
matching determines the boundary data of its new harmonic term from
the degree $2r+j$ Taylor term and all preceding coefficients. Harmonic
extension determines this term. Derivative matching is then
$4r^2b_j+\text{known forcing}=0$. The new $b_j$ does not enter the
order-$j$ value equation. Nonlinear lower-order terms contribute when
$j\geq2$.

The truncated curve changes the physical center by $O(R^{J+2})$.
Using it in the geometric correction is valid whenever
\[
 \sqrt q R^{J+1}\to0,
 \quad\hbox{equivalently}\quad
 q=o(m^{(J+1)/(r+J+1)}).
\]
If $b_1=\cdots=b_{j-1}=0$ but $b_j\not\equiv0$, the exact criterion
\eqref{eq:fixed-weight-jet-criterion} makes the fixed-weight uncorrected range
$q=o(m^{j/(r+j)})$, sharply in the all-direction sense. Higher terms
are relatively $O(R)$ at an angle where $b_j\ne0$.

In the radial case $\phi=f(u)$, $u=|z|^2$, the hole is a disk by
uniqueness and rotational invariance. Its exact center satisfies
\[
 m u_*f'(u_*)=q.
\]
The asymptotic mode variance is $m u\rho(u)$, with
$\rho(u)=f'(u)+uf''(u)$. The corresponding Gaussian argument is
\begin{equation}\label{eq:radial-growing-jet-argument}
 \frac{m u f'(u)-q}{\sqrt{m u\rho(u)}}.
\end{equation}
For bounded arguments, Taylor expansion at $u_*$ identifies it with
$2\sqrt{m\rho(u_*)}(\sqrt u-\sqrt{u_*})+o(1)$, so
\Cref{thm:curvature-boundary-jets} proves this radial formula for all sublinear
orders. For $f=u^r+c u^{r+1}+\cdots$, its first raw-center shift is
$c(r+1)r^{-3/2}\sqrt q R^2$, with critical growth $m^{2/(r+2)}$.
This larger radial threshold does not describe arbitrary coordinates.

No fixed finite Taylor jet determines the correction for every
$q=o(m)$. For example, near zero take
\[
 f(u)=u^r+c u^r \mathrm{e}^{-1/\sqrt u},\qquad c>0,
\]
smoothly extended at zero and interpolated globally as above.
It has the same Taylor jet as $u^r$ at every finite order, but at
$u=R^2$ its relative contribution to $uf'(u)$ is
$c \mathrm{e}^{-1/R}(1+1/(2rR))$. For $q=\lfloor m/\log m\rfloor$,
$\sqrt q \mathrm{e}^{-1/R}/R\to\infty$. The raw observation with $\tau=0$
is outside the true center by many transition widths, giving limit one.
The exact quadrature-domain correction includes this flat contribution.
